\chapter{Fundamentação Teórica}
\label{sec:fundamentacao}

Este capítulo apresenta os conceitos que sustentam a análise da unificação de sistemas proposta neste trabalho. O fluxo segue do geral para o específico: inicia-se com os conceitos de Sistemas de Informação e a forma como são avaliados, em seguida, é discutido os sistemas legados e os efeitos da fragmentação. depois, são apresentadas as abordagens de integração e de unificação, com seus benefícios e riscos. Na sequência, iremos tratar a gestão de processos de negócio e a aceitação de tecnologia pelos usuários, fatores que determinam se uma solução unificada gera, de fato, resultados. Por fim, contextualizaremos o setor elétrico brasileiro, em especial o Ambiente de Contratação Regulada, no qual se insere o caso estudado. A descrição do aplicativo InLight e das tecnologias empregadas em sua construção é apresentada no Capítulo~\ref{sec:proposta}.

\section{Sistemas de Informação nas Organizações}
\label{sec:si}

\subsection{Conceito e Componentes}

Sistemas de Informação (SI) são conjuntos de componentes inter-relacionados que coletam, processam, armazenam e distribuem informações para apoiar a tomada de decisão, a coordenação e o controle em uma organização \citep{laudon2020}. Em definição semelhante, \citet{obrien2013} descrevem o SI como uma combinação organizada de pessoas, hardware, software, redes de comunicação, recursos de dados, políticas e procedimentos que armazenam, recuperam, transformam e disseminam informações.

Os componentes básicos de um SI são \citep{stair2018}:

\begin{itemize}
    \item \textbf{Hardware:} equipamentos utilizados na entrada, no processamento e na saída de dados;
    \item \textbf{Software:} programas que controlam o hardware e executam as tarefas do sistema;
    \item \textbf{Dados:} fatos brutos que, processados, tornam-se informação útil;
    \item \textbf{Redes:} infraestrutura de comunicação que conecta os demais componentes;
    \item \textbf{Pessoas:} usuários, desenvolvedores e gestores que interagem com o sistema;
    \item \textbf{Procedimentos:} políticas e regras que orientam a operação do sistema.
\end{itemize}

Essa visão evidencia que um SI não se resume à tecnologia: ele envolve pessoas e procedimentos, o que será retomado ao se discutir a adoção de sistemas unificados (Seção~\ref{sec:aceitacao}). Para \citet{turban2010}, os SI são essenciais às organizações modernas porque transformam dados brutos em informações que sustentam decisões estratégicas, táticas e operacionais. A qualidade dessas decisões depende diretamente da qualidade e da disponibilidade das informações fornecidas.

\subsection{Avaliação do Sucesso de Sistemas de Informação}

Para analisar os impactos de um novo sistema, é necessário definir o que se entende por ``sucesso''. O modelo de \citet{delone2003}, um dos mais utilizados na área, organiza o sucesso de um SI em seis dimensões inter-relacionadas:

\begin{itemize}
    \item \textbf{Qualidade do sistema:} características técnicas, como facilidade de uso, confiabilidade e tempo de resposta;
    \item \textbf{Qualidade da informação:} completude, precisão, atualidade e relevância das informações geradas;
    \item \textbf{Qualidade do serviço:} suporte oferecido aos usuários pela equipe responsável pelo sistema;
    \item \textbf{Uso (ou intenção de uso):} grau em que o sistema é efetivamente utilizado;
    \item \textbf{Satisfação do usuário:} percepção dos usuários em relação ao sistema;
    \item \textbf{Benefícios líquidos:} impactos positivos e negativos do sistema sobre indivíduos, grupos e a organização.
\end{itemize}

Esse modelo oferece uma estrutura para observar e mensurar os benefícios e as limitações da unificação de sistemas, uma das questões de pesquisa deste trabalho, e será usado como referência na análise dos resultados.

\section{Sistemas Legados e Fragmentação}
\label{sec:legados}

\subsection{Sistemas Legados}

\citet{bennett1995} caracteriza os sistemas legados como sistemas de software de grande porte que continuam vitais para a organização, mas que ela não sabe mais como manter ou evoluir adequadamente. \citet{bisbal1999} os definem como sistemas de informação que resistem de forma significativa a modificações e à evolução. Em geral, foram desenvolvidos com tecnologias que se tornaram ultrapassadas, têm documentação escassa ou desatualizada e concentram regras de negócio críticas embutidas no código, o que torna sua manutenção cada vez mais cara e arriscada \citep{bisbal1999,pressman2016}.

\citet{bisbal1999} observam que esses sistemas criam um dilema: substituí-los envolve alto risco, mas mantê-los limita a capacidade de a organização responder a novas necessidades. Os autores agrupam as alternativas em três estratégias: o \textbf{redesenvolvimento}, que reconstrói o sistema do zero; o \textbf{encapsulamento} (\textit{wrapping}), que oferece uma nova interface sobre o sistema existente; e a \textbf{migração}, que transfere o sistema para um ambiente mais flexível preservando seus dados e funcionalidades. Essa classificação é útil para situar o InLight, que não substitui os sistemas de origem, mas consolida suas informações em uma nova aplicação.

\subsection{Fragmentação de Sistemas e seus Impactos}

A fragmentação ocorre quando uma organização utiliza diversas aplicações isoladas, cada uma responsável por uma parte dos processos de negócio, sem integração adequada entre elas \citep{laudon2020}. Esse cenário é comum em empresas que adotaram novas tecnologias de forma incremental, sem uma arquitetura que as conectasse, e produz os efeitos descritos a seguir.

\subsubsection{Redundância e Inconsistência de Dados}

Quando a mesma informação é armazenada em vários sistemas, cada um com seus próprios arquivos, surgem a redundância e o risco de inconsistência, pois uma atualização feita em um local pode não ser refletida nos demais \citep{elmasri2016}. Na prática, dados de um mesmo contrato podem existir em versões diferentes em sistemas distintos e em planilhas, gerando dúvidas sobre qual delas é a correta.

\subsubsection{Retrabalho e Perda de Produtividade}

A ausência de integração obriga os colaboradores a executar tarefas manuais repetitivas, como digitar os mesmos dados em mais de um sistema, exportar e importar planilhas e conciliar informações manualmente \citep{stair2018}. Além de consumir tempo, cada transcrição manual é uma oportunidade para erros.

\subsubsection{Dificuldade na Tomada de Decisão}

A fragmentação impede uma visão consolidada dos processos. Os gestores precisam reunir manualmente informações de diversas fontes antes de analisá-las, um processo demorado e sujeito a erros \citep{power2002}. Para \citet{sharda2014}, a integração de dados é condição para que sistemas de apoio à decisão sejam eficazes.

\subsubsection{Custos Operacionais e Riscos de Segurança}

Manter vários sistemas não integrados implica custos de infraestrutura, licenciamento, treinamento e suporte \citep{turban2010}. Além disso, ambientes fragmentados dificultam a aplicação de políticas de segurança e de controle de acesso uniformes, já que cada sistema pode ter mecanismos e níveis de proteção distintos \citep{whitman2018}.

\subsection{Custos e Benefícios da Integração de Dados}

Apesar dos problemas apontados, a integração não deve ser tratada como benéfica em qualquer circunstância. \citet{goodhue1992} questionam a suposição de que integrar dados sempre gera ganhos líquidos: a integração melhora a coordenação e a tomada de decisão em toda a organização, mas pode reduzir a autonomia e a flexibilidade das áreas locais, além de alterar os custos de projeto e implantação dos sistemas. Para esses autores, os benefícios superam os custos apenas em determinadas situações, e não necessariamente para todos os dados da organização. Essa ressalva orienta a análise deste trabalho: o objetivo não é presumir que a unificação é sempre vantajosa, mas identificar em quais processos ela gera valor.

\section{Integração e Unificação de Sistemas}
\label{sec:integracao}

\subsection{Integração de Sistemas}

A integração de sistemas consiste em conectar aplicações distintas para que compartilhem dados e processos de negócio \citep{hohpe2003}. Segundo \citet{linthicum2000}, integrar não significa necessariamente substituir os sistemas existentes, mas criar mecanismos que permitam a comunicação entre eles. O autor distingue níveis de integração conforme o ponto em que as aplicações se conectam:

\begin{itemize}
    \item \textbf{Nível de dados:} movimentação de dados diretamente entre os bancos de dados das aplicações;
    \item \textbf{Nível de interface de aplicação:} uso das interfaces de programação (APIs) que as aplicações expõem;
    \item \textbf{Nível de método:} compartilhamento da lógica de negócio entre aplicações;
    \item \textbf{Nível de interface de usuário:} uso da própria interface do usuário como ponto de integração, geralmente em sistemas legados que não oferecem outro acesso.
\end{itemize}

\citet{hohpe2003} descrevem quatro estilos básicos de integração entre aplicações: \textbf{transferência de arquivos}, em que uma aplicação gera arquivos que outra consome; \textbf{banco de dados compartilhado}, em que várias aplicações leem e gravam na mesma base; \textbf{invocação remota de procedimentos}, em que uma aplicação chama funcionalidades de outra; e \textbf{mensageria}, em que as aplicações trocam mensagens de forma assíncrona. Sobre esses estilos foram construídas tecnologias como as APIs REST \citep{fielding2000}, os \textit{Web Services} baseados em SOAP \citep{alonso2004}, os barramentos de serviço corporativos (ESB) \citep{chappell2004} e as arquiteturas de microsserviços \citep{newman2015}.

\subsection{Unificação de Sistemas}

Enquanto a integração conecta sistemas que continuam separados, a unificação consolida funcionalidades e informações dispersas em uma plataforma única ou em um conjunto coeso de sistemas \citep{laudon2020}. O exemplo mais conhecido são os sistemas integrados de gestão (ERP), que, segundo \citet{davenport1998}, prometem integrar as informações que circulam pela empresa por meio de uma base de dados comum, atualizada em tempo real e acessada por diferentes áreas.

Em comparação com a integração pontual, a unificação oferece:

\begin{itemize}
    \item \textbf{Interface única e consistente:} o usuário passa a trabalhar em uma só aplicação, sem alternar entre sistemas com lógicas diferentes; a consistência da interface é um dos princípios centrais de usabilidade \citep{nielsen1993,norman2013};
    \item \textbf{Base de dados comum:} os dados passam a ter uma origem única, o que reduz a redundância e os conflitos de versão \citep{elmasri2016};
    \item \textbf{Manutenção concentrada:} melhorias e correções são aplicadas em um só lugar \citep{pressman2016};
    \item \textbf{Base para a evolução:} plataformas modulares facilitam a inclusão de novas funcionalidades e a adaptação a mudanças no negócio \citep{ross2006}.
\end{itemize}

\subsection{Benefícios Esperados}

Para classificar os benefícios de sistemas corporativos integrados, \citet{shang2002} propõem um modelo com cinco dimensões, construído a partir de casos reais de implantação:

\begin{itemize}
    \item \textbf{Operacionais:} redução de custos e de tempo de ciclo, aumento de produtividade e melhoria da qualidade;
    \item \textbf{Gerenciais:} melhor gestão de recursos, melhor tomada de decisão e melhor controle do desempenho;
    \item \textbf{Estratégicos:} apoio ao crescimento do negócio e a novas iniciativas;
    \item \textbf{De infraestrutura de TI:} maior flexibilidade, redução de custos de TI e ampliação da capacidade da infraestrutura;
    \item \textbf{Organizacionais:} mudanças nos padrões de trabalho, aprendizado organizacional e foco comum entre as áreas.
\end{itemize}

Esse modelo complementa o de \citet{delone2003}: enquanto este define as dimensões do sucesso de um SI, aquele detalha os tipos de benefício que se pode esperar de um sistema integrado. Ambos serão utilizados para organizar a análise dos impactos do InLight.

\subsection{Desafios e Fatores Críticos de Sucesso}

A unificação também envolve riscos. \citet{davenport1998} observa que muitos projetos de sistemas integrados fracassam por razões de negócio, e não técnicas: as organizações não conciliam as exigências da tecnologia com suas próprias necessidades. Entre os principais desafios estão:

\begin{itemize}
    \item \textbf{Complexidade técnica:} migração de dados legados, mapeamento de regras de negócio e compatibilidade com sistemas existentes \citep{bisbal1999};
    \item \textbf{Resistência à mudança:} usuários habituados aos sistemas antigos tendem a resistir a novos modos de trabalho \citep{kotter1996};
    \item \textbf{Custos e riscos de transição:} investimento inicial e possibilidade de falhas durante a implantação \citep{laudon2020};
    \item \textbf{Perda de autonomia local:} áreas acostumadas a controlar os próprios dados podem perder flexibilidade \citep{goodhue1992}.
\end{itemize}

\citet{nah2001} identificaram onze fatores críticos para o sucesso na implantação de sistemas integrados, entre eles: apoio da alta direção, composição da equipe do projeto, programa de gestão da mudança, plano e visão de negócio, redesenho de processos com mínima customização, gestão do projeto, monitoramento e avaliação de desempenho, comunicação eficaz e tratamento adequado dos sistemas legados. Esses fatores mostram que o resultado depende tanto da tecnologia quanto da forma como a organização conduz a mudança.

\section{Processos de Negócio e Automação}
\label{sec:processos}

\subsection{Processos de Negócio e sua Gestão}

Um processo de negócio é um conjunto estruturado de atividades, com início e fim definidos, que produz um resultado de valor para um cliente ou mercado \citep{davenport1993}. A Gestão de Processos de Negócio (BPM, do inglês \textit{Business Process Management}) é a disciplina que supervisiona como o trabalho é executado para garantir resultados consistentes e aproveitar oportunidades de melhoria \citep{dumas2018}. \citet{dumas2018} organizam essa disciplina em um ciclo de vida composto pelas etapas de identificação, descoberta (modelagem do processo atual), análise, redesenho, implementação e monitoramento dos processos.

Para representar processos, a notação mais difundida é a BPMN (\textit{Business Process Model and Notation}), padrão mantido pelo \textit{Object Management Group} \citep{dumas2018}. A modelagem permite identificar ineficiências, gargalos, redundâncias e oportunidades de melhoria antes de implementar mudanças \citep{aguilarsaven2004,sharp2009}.

\subsection{Reengenharia: Redesenhar antes de Automatizar}

\citet{hammerchampy1993} definem a reengenharia como o repensar fundamental e o redesenho radical dos processos de negócio para alcançar melhorias expressivas em medidas críticas de desempenho, como custo, qualidade, serviço e velocidade. Antes disso, \citet{hammer1990} já alertava que usar a tecnologia apenas para acelerar processos existentes não resolve suas deficiências; o ganho surge quando os processos são repensados a partir das possibilidades da tecnologia. Entre os princípios propostos pelo autor está o de capturar a informação uma única vez, na fonte, eliminando a reentrada de dados.

Esse princípio é diretamente relevante para a unificação de sistemas: transferir para um novo aplicativo os mesmos passos manuais de antes traria ganhos limitados. O benefício depende de eliminar etapas redundantes, como a digitação repetida de dados em sistemas e planilhas diferentes.

\subsection{Automação de Processos e Workflows}

Um \textit{workflow} pode ser entendido como um conjunto de tarefas organizadas para realizar um processo de negócio, executadas por pessoas ou por sistemas \citep{georgakopoulos1995}. Sistemas de \textit{workflow} coordenam a sequência de atividades, atribuem tarefas, controlam prazos e registram o histórico de execução \citep{vanderaalst2003}.

Uma forma mais recente de automação é a \textit{Robotic Process Automation} (RPA), em que robôs de software reproduzem as ações que um usuário executaria na interface dos sistemas \citep{vanderaalst2018}. A RPA é indicada para tarefas repetitivas e de alto volume e, por não exigir mudanças nos sistemas existentes, é uma alternativa à integração em ambientes com sistemas legados \citep{lacity2016,vanderaalst2018}. Ela não elimina, porém, a fragmentação dos dados, apenas automatiza a passagem de informações entre sistemas que continuam separados.

\subsection{Medição de Desempenho}

Para verificar se uma mudança gerou melhoria, é preciso medir o desempenho dos processos. \citet{dumas2018} propõem quatro dimensões de desempenho: \textbf{tempo} (por exemplo, o tempo de ciclo do processo), \textbf{custo}, \textbf{qualidade} (incluindo a taxa de erros e de retrabalho) e \textbf{flexibilidade}. Os autores ressaltam que essas dimensões costumam competir entre si: melhorar uma delas pode piorar outra. Em nível organizacional, o \textit{Balanced Scorecard} de \citet{kaplan1996} recomenda avaliar o desempenho sob quatro perspectivas, financeira, de clientes, de processos internos e de aprendizado e crescimento, evitando uma visão limitada a indicadores financeiros.

\section{Aceitação de Tecnologia pelos Usuários}
\label{sec:aceitacao}

Um sistema só gera benefícios se for efetivamente utilizado. O Modelo de Aceitação de Tecnologia (TAM, do inglês \textit{Technology Acceptance Model}), proposto por \citet{davis1989}, explica a adoção de uma tecnologia por dois fatores principais: a \textbf{utilidade percebida}, que é o grau em que a pessoa acredita que o sistema melhora seu desempenho no trabalho, e a \textbf{facilidade de uso percebida}, que é o grau em que acredita que usá-lo exigirá pouco esforço. Os estudos do autor indicaram que a utilidade percebida tem relação mais forte com o uso do sistema e que a facilidade de uso atua, em grande parte, aumentando a utilidade percebida.

\citet{venkatesh2003} ampliaram essa abordagem com a Teoria Unificada de Aceitação e Uso de Tecnologia (UTAUT), que reúne oito modelos anteriores em quatro determinantes:

\begin{itemize}
    \item \textbf{Expectativa de desempenho:} crença de que o sistema ajudará a obter ganhos no trabalho;
    \item \textbf{Expectativa de esforço:} facilidade associada ao uso do sistema;
    \item \textbf{Influência social:} percepção de que pessoas importantes, como gestores e colegas, esperam que o sistema seja usado;
    \item \textbf{Condições facilitadoras:} existência de infraestrutura técnica e organizacional para apoiar o uso.
\end{itemize}

Esses modelos fundamentam a questão de pesquisa sobre como a visão de usuários e gestores influencia a adoção de sistemas unificados. No caso do InLight, indicam que a adoção dependerá de os usuários do BackOffice perceberem ganhos concretos em relação às planilhas e aos sistemas atuais, e do apoio dos gestores à mudança.

\section{O Setor Elétrico Brasileiro e a Comercialização de Energia}
\label{sec:setor}

\subsection{Ambientes de Contratação: ACR e ACL}

O modelo vigente de comercialização de energia elétrica no Brasil foi instituído pela Lei nº 10.848/2004 \citep{brasil2004lei} e regulamentado pelo Decreto nº 5.163/2004 \citep{brasil2004decreto}. Esse marco divide a comercialização em dois ambientes:

\begin{itemize}
    \item \textbf{Ambiente de Contratação Regulada (ACR):} reúne as compras de energia das distribuidoras para atender seus consumidores cativos. As contratações são feitas, em regra, por meio de leilões regulados pela ANEEL e operacionalizados pela CCEE \citep{brasil2004lei};
    \item \textbf{Ambiente de Contratação Livre (ACL):} reúne as operações de compra e venda de energia por meio de contratos bilaterais livremente negociados entre geradores, comercializadores e consumidores livres \citep{brasil2004decreto}.
\end{itemize}

A mesma lei autorizou a criação da Câmara de Comercialização de Energia Elétrica (CCEE), responsável, entre outras atribuições, pelo registro dos contratos, pela contabilização da energia e pela liquidação financeira das operações no mercado de curto prazo \citep{brasil2004lei,ccee2023}.

\subsection{Contratos do Ambiente Regulado}

No ACR, os contratos firmados entre vendedores e distribuidoras a partir dos leilões são denominados Contratos de Comercialização de Energia Elétrica no Ambiente Regulado (CCEAR). O Decreto nº 5.163/2004 prevê duas modalidades \citep{brasil2004decreto}:

\begin{itemize}
    \item \textbf{CCEAR por quantidade de energia:} o vendedor se compromete a entregar uma quantidade de energia, assumindo os riscos hidrológicos;
    \item \textbf{CCEAR por disponibilidade de energia:} o vendedor disponibiliza a capacidade de geração, e os riscos hidrológicos são assumidos pelo comprador, com repasse aos consumidores.
\end{itemize}

Além dessas modalidades, a Lei nº 12.783/2013 estabeleceu o regime de \textbf{cotas de garantia física}, pelo qual a energia de usinas hidrelétricas com concessões prorrogadas é alocada às distribuidoras na forma de cotas definidas pela ANEEL \citep{brasil2013lei}. Contratos de quantidade, de disponibilidade e de cotas correspondem às categorias tratadas pelo aplicativo estudado neste trabalho (Seção~\ref{sec:inlight}).

Cada um desses contratos envolve dados de origens diferentes, como os registros da CCEE, os sistemas internos da distribuidora e controles paralelos em planilhas, além de prazos regulatórios para registro, contabilização e liquidação. Esse cenário reúne as condições descritas na Seção~\ref{sec:legados}: informações dispersas, conciliação manual e risco de inconsistência.

\subsection{Transformação Digital no Setor Elétrico}

O setor elétrico passa por uma transformação impulsionada pelas redes inteligentes (\textit{smart grids}), pela medição digital e pelo uso intensivo de dados. \citet{farhangi2010} aponta que o aumento do custo da energia, a eletrificação crescente e as mudanças climáticas são os principais vetores dessa transformação, que exige das empresas sistemas capazes de tratar volumes crescentes de informação. Para \citet{westerman2014}, porém, a transformação digital não se limita à tecnologia: ela exige repensar processos e modelos de negócio, o que reforça a relação entre unificação de sistemas e redesenho de processos discutida na Seção~\ref{sec:processos}.

\section{Considerações sobre o Capítulo}

Os conceitos apresentados formam a base para a análise dos capítulos seguintes. A literatura mostra que a fragmentação de sistemas gera redundância, retrabalho e dificuldade de decisão \citep{laudon2020,elmasri2016}, mas também que a integração tem custos e não é benéfica em todas as situações \citep{goodhue1992}. Os benefícios da unificação podem ser classificados com o modelo de \citet{shang2002} e avaliados com as dimensões de \citet{delone2003}; sua concretização depende do redesenho dos processos \citep{hammer1990,dumas2018}, da gestão da mudança \citep{nah2001} e da aceitação pelos usuários \citep{davis1989,venkatesh2003}. No setor elétrico, a variedade de contratos do ACR e a dependência de informações externas, como as da CCEE, tornam esses fatores especialmente relevantes. O Capítulo~\ref{sec:revisao} examina como estudos empíricos trataram essas questões, e o Capítulo~\ref{sec:proposta} aplica esses conceitos ao caso do aplicativo InLight.
